% !TeX program = xelatex
\documentclass[UTF8,a4paper,12pt]{ctexart}

\usepackage{array}
\usepackage{booktabs}
\usepackage{geometry}
\usepackage{longtable}
\usepackage{setspace}
\usepackage{tabularx}

\setmainfont{Times New Roman}
\setCJKmainfont{SimSun}[AutoFakeBold=2.5]
\setCJKsansfont{SimHei}

\geometry{
  a4paper,
  top=2.54cm,
  bottom=2.54cm,
  left=2.70cm,
  right=2.70cm
}

\setlength{\parindent}{2em}
\linespread{1.25}
\raggedbottom

\title{\heiti AI 工具使用说明}
\author{}
\date{}

\begin{document}
\maketitle

\section{所用 AI 工具名称和版本}

\begin{table}[h]
\centering
\caption{AI 工具使用情况}
\begin{tabularx}{\textwidth}{>{\centering\arraybackslash}p{2.8cm}>{\centering\arraybackslash}p{3.6cm}>{\centering\arraybackslash}X>{\centering\arraybackslash}p{4.0cm}}
\toprule
工具名称 & 版本/型号 & 开发机构/公司 & 使用日期\\
\midrule
DeepSeek & DeepSeek-V2.5-Pro & 深度求索（DeepSeek） & 2026-06-19 至 2026-06-22\\
MiMo & MiMo & 小米（Xiaomi） & 2026-06-19 至 2026-06-22\\
\bottomrule
\end{tabularx}
\end{table}

\section{具体使用目的和环节}

本论文写作过程中，AI 工具主要用于辅助文本组织、语言润色、格式规范和 LaTeX 排版检查。具体使用环节如下：

\begin{enumerate}
  \item 论文结构整合：辅助梳理问题重述、问题分析、模型假设、符号说明、模型建立与求解、模型评价和参考文献之间的行文顺序。
  \item 摘要与关键词润色：根据论文实际求解内容，辅助润色摘要表述，使其覆盖四个问题的建模思路、求解结果和模型特点。
  \item LaTeX 语法整理：辅助检查公式、表格、图片引用、脚注、目录和参考文献格式，使论文排版更加规范。
  \item 引用与参考文献格式检查：辅助在正文相应位置加入文献引用，并将 AI 工具使用情况补充到参考文献中。
\end{enumerate}

\section{关键交互记录}

以下记录为关键交互的摘要式摘录，保留重要提示词、回复要点和人工处理情况，不列出完整对话。

\small
\begin{longtable}{>{\centering\arraybackslash}p{0.8cm}p{3.5cm}p{4.8cm}p{3.6cm}}
\caption{关键交互记录摘要}\\
\toprule
序号 & 重要提示词 & AI 回复要点 & 采纳与人工处理\\
\midrule
\endfirsthead
\toprule
序号 & 重要提示词 & AI 回复要点 & 采纳与人工处理\\
\midrule
\endhead
\bottomrule
\endfoot
1 & 将问题一、二、三内容整合进最终论文，保留图表和关键结果。 & 建议按 5.1、5.2、5.3 分节组织，复制图片到统一目录，并保持核心公式和数值结果。 & 采纳章节组织方式，人工核对图片、表格和公式内容。\\
2 & 前三问模型建立与求解存在重复，进行中等压缩。 & 建议保留主模型、关键公式、图表和数值结果，压缩重复推导和过程性叙述。 & 采纳文字压缩方向，人工确认频率、振幅、初相和残差指标未被改变。\\
3 & 根据 Word 模板补入第四问和模型评价。 & 建议将问题四按传感器布局检测模型、AHP 参数确定、布局优化和验证分析组织。 & 采纳结构整理，人工检查最优布局、检测概率和敏感性分析表格。\\
4 & 修改符号说明。 & 将重要符号汇总为三列表格。 & 采纳格式建议，人工确认符号取舍。\\
\end{longtable}
\normalsize

\section{采纳和人工修改情况}

AI 工具提供的内容主要作为论文写作和排版的辅助材料。论文中的模型选择、求解流程、数值结果、图表内容和最终结论均由参赛队员人工撰写、绘制。

在采纳AI回答过程中，人工主要进行了以下修改：
\begin{enumerate}
  \item 对 AI 给出的语言润色内容进行筛选，删除与题目无关或过度展开的表述。
  \item 对所有核心公式、频率估计值、振幅、初相、残差指标和检测概率结果进行人工核对。
  \item 对图表标题、符号说明、模型假设和参考文献格式进行人工调整，使其符合数学建模论文写作习惯。
  \item 对 LaTeX 编译结果进行检查，确保主论文和本说明文件能够正常生成 PDF。
\end{enumerate}

综上，AI 工具在本论文中承担辅助润色、结构建议、LaTeX 语法整理和格式检查作用；论文的核心建模思路、结果解释和最终提交内容由人工负责。

\end{document}
